\section{Positive traces and a quantitative charge estimate}
\label{sec:transfer}

The finite-dimensional argument allows spatial coefficients and
the mass matrix to vary in time, as required for quantized link
histories. In this section the one-particle dimension is $D=2d$,
with the same four-component Dirac structure as above.

\begin{lemma}[Positivity and the transfer annulus]
\label{lem:transfer}
Consider ordered free evolutions generated by
\[
 h(t)=\sum_j\alpha_jv_j(t)+\beta M(t),\qquad M(t)\ge m_*I>0,
\]
with bounded Hermitian spatial coefficients and total free time
$b>0$, interspersed with spin-scalar unitary kicks.
Let $U$ be their one-particle product. Include each kick's
centering phase in the finite-CAR trace
\begin{equation}
 F=e^{\frac12\int\Tr A(t)\,dt}\det(I+U),
 \qquad U'=-A(t)U,\quad A=c(t)h(t)+i\Phi(t).
 \label{eq:centered-trace}
\end{equation}
The auxiliary ODE interval is $[0,T]$; free and kick intervals
can be separate, $c\ge0$ obeys $\int_0^Tc(t)\,dt=b$, and
$\Phi$ is Hermitian and spin scalar.
Then $F>0$. Exactly $d$ transfer eigenvalues, with algebraic
multiplicity, obey
\begin{equation}
 |\lambda|\le e^{-bm_*};
 \label{eq:annulus}
\end{equation}
the remaining $d$ are their partners $1/\overline\lambda$ and
have modulus at least $e^{bm_*}$. Normality or diagonalizability
of $U$ is not required.
\end{lemma}
\begin{proof}
Since $C_5$ anticommutes with $h$ and commutes with $\Phi$,
\[
 A^*C_5+C_5A=0,\qquad U^*C_5U=C_5.
\]
Thus $U^*=C_5U^{-1}C_5$ and the spectrum is paired, with
algebraic multiplicities, under $\lambda\mapsto1/\overline\lambda$.
Also
\[
 \overline{\det(I+U)}=\det(U)^{-1}\det(I+U),\qquad
 \det U=e^{-\int\Tr A}.
\]
Because $\Tr h=0$, $\int\Tr A$ is purely imaginary. These
identities make \eqref{eq:centered-trace} real.

For any transfer eigenvector $v$, transport it along the
successive intervals and kicks, obtaining $f(t)$.
The form $B(t)=\langle f(t),\beta f(t)\rangle$ is unchanged by
scalar kicks and satisfies
\begin{equation}
 B'(t)=-2c(t)\langle f(t),M(t)f(t)\rangle
       \le-2m_*c(t)\|f(t)\|^2.
 \label{eq:beta-form}
\end{equation}
Its endpoint is $B(T)=|\lambda|^2B(0)$, and its total decrease
is strict because $b>0$ and a nonzero solution never vanishes.
Hence $|\lambda|=1$ is impossible.
If $|\lambda|<1$, both endpoints and the entire trajectory
have $B>0$. Since $\|f\|^2\ge B$, integration of
$B'\le-2m_*c(t)B$ proves \eqref{eq:annulus}.
If $|\lambda|>1$, the entire trajectory has $B<0$;
$\|f\|^2\ge-B$ gives $B'\le2m_*c(t)B$ and
$|\lambda|\ge e^{bm_*}$, reversing the inequality when
dividing by $B(0)<0$. Spectral pairing gives exactly $d$
eigenvalues in each region.

In particular $\det(I+U)\ne0$. Continuously turn off all kicks
and all $v_j$, and replace $M(t)$ by
$(1-s)M(t)+sm_*I$. The real centered trace never vanishes
along this deformation. At its endpoint it equals
$\det(I+e^{-bm_*\beta})>0$, which proves its sign.
The same argument holds for bounded measurable histories:
their solutions are absolutely continuous and
\eqref{eq:beta-form} holds almost everywhere.
\end{proof}

For a scalar kick $u=e^{i\Theta}$ the centering factor is
$e^{-\frac i2\Tr\Theta}$. For a local gauge twist,
$\Tr\Theta=4\sum_x\theta_x$, so it is precisely
$e^{-2i\sum_x\theta_x}$ from $q_x=n_x-2$.
This phase is indispensable and is single-valued when
$\theta_x$ changes by $2\pi$.
For a nonnegative mass coefficient, replace $M$ by
$M+\epsilon I$ and pass to the finite-dimensional limit.
The centered trace is then nonnegative but may vanish.

\begin{lemma}[Canonical Fourier bound]
\label{lem:fourier}
Append to a history of Lemma~\ref{lem:transfer} a global
charge twist $e^{i\phi Q}$, $Q=d\Gamma(I)-d$.
Let $F(\phi)$ be its centered trace and
$F_q=(2\pi)^{-1}\int_0^{2\pi}e^{-iq\phi}F(\phi)\,d\phi$.
For $|q|\le d$,
\begin{equation}
 |F_q|\le\binom{2d}{d+|q|}e^{-bm_*|q|}F_0.
 \label{eq:fourier-bound}
\end{equation}
\end{lemma}
\begin{proof}
Let $\lambda_1,\ldots,\lambda_d$ be the transfer eigenvalues
in the small disk. Pairing determinant factors and including
the phase $e^{-id\phi}$ gives
\[
 F(\phi)=C\prod_{j=1}^d|1+e^{i\phi}\lambda_j|^2,\qquad C>0.
\]
A paired factor is
$e^{i\phi}\overline\lambda_j^{-1}
 |1+e^{i\phi}\lambda_j|^2$; the remaining constant is positive
because $F(0)>0$. This uses determinant eigenvalues only and
also holds for defective matrices.
If $e_j$ are the elementary symmetric polynomials in the
$\lambda_j$ and $s=e^{-bm_*}$, then
\[
 F_0=C\sum_{j=0}^d|e_j|^2\ge C,\qquad
 |e_j|\le\binom dj s^j.
\]
For $q\ge0$, Vandermonde's convolution yields
\[
 |F_q|\le Cs^q\sum_{j=0}^{d-q}\binom d{j+q}\binom dj s^{2j}
 \le Cs^q\binom{2d}{d+q}
 \le s^q\binom{2d}{d+q}F_0.
\]
Negative $q$ follows by complex conjugation.
\end{proof}

The prefactor in \eqref{eq:fourier-bound} need not be one.
It is retained until division by inverse temperature in the
ground-energy limit.
